\documentclass[conference]{IEEEtran}
\IEEEoverridecommandlockouts

\usepackage{cite}
\usepackage{amsmath,amssymb,amsfonts}
\usepackage{algorithm}
\usepackage{algorithmic}
\usepackage{graphicx}
\usepackage{textcomp}
\usepackage{xcolor}
\usepackage{url}
\usepackage{hyperref}

\hypersetup{hidelinks}

\title{Resilience-Aware, Human-Centric AI-RAN Orchestration for Service Continuity Across Terrestrial and Non-Terrestrial Networks}

\author{\IEEEauthorblockN{Edmund Gunn Jr.}
\IEEEauthorblockA{\textit{7GC Research Product Spine}\\
Email: CITATION\_NEEDED}}

\begin{document}

\maketitle

\begin{abstract}
We describe a methods-ready research artifact for privacy-preserving, phone-first field measurement and twin-informed AI-RAN orchestration under heterogeneous terrestrial and non-terrestrial connectivity.
The artifact couples canonical JSON Schema contracts, an integrated Gate~2 pipeline, preregistered evaluation design, and a 54-cell controlled pilot matrix (3 zones $\times$ 2 network conditions $\times$ 3 days $\times$ 3 workload profiles).
We explicitly disclaim carrier-grade AI-RAN, deployable 6G infrastructure, citywide generalization, and unauthorized RF transmission.
Primary outcome \texttt{total\_service\_outage\_time\_s} (secondary \texttt{time\_to\_recovery\_s}; Amendment~001 in progress) is defined from physical probes; empirical pilot results are \textbf{RESULTS\_PENDING\_AUTHENTIC\_GATE3\_DATA} with no outcome numbers reported here.
\end{abstract}

\begin{IEEEkeywords}
6G-aligned, IMT-2030, AI-RAN, O-RAN, digital twin, non-terrestrial networks, reproducibility, privacy-preserving measurement
\end{IEEEkeywords}

\section{Introduction}
\label{sec:introduction}

Sparse, degraded, and disrupted connectivity remains a practical barrier to service continuity for learning, creation, and sensing workloads in under-connected communities.
IMT-2030-aligned research explores AI-native air interfaces, open RAN disaggregation, and multi-access resilience including non-terrestrial network (NTN) fallback~\cite{itu_m2160,saad2020vision,strinati20216g}.
Operational claims about carrier-grade AI-RAN or deployable 6G infrastructure are outside the scope of this artifact.

This paper documents the \emph{methods-ready} umbrella repository \texttt{gunnchos-7gc-ai-ran-field-kit}, which locks a doctoral thesis hypothesis and three research questions (RQ1--RQ3) in Gate~1, provides executable contracts and orchestration in Gate~2, defines controlled evidence collection in Gate~3, and preregisters scientific evaluation in Gate~4.
We separate engineering completion from release evidence: Gate~2 system pass does not imply immutable archive, DOI deposit, or non-author reproduction.

\subsection{Central hypothesis}
A privacy-preserving orchestration architecture that combines physical device measurements, workload and service intent, digital-twin context, and terrestrial, local-edge, and NTN state can improve service continuity and worst-user quality of experience under sparse connectivity when compared with static, network-only, and service-priority baselines, while maintaining explicit fairness, energy, privacy, and reliability constraints.

\subsection{Research questions}
\textbf{RQ1 (Trustworthy context):} How can privacy-preserving Edge-IO measurements be transformed into reproducible digital-twin state without direct identifiers?
\textbf{RQ2 (Human-centric orchestration):} To what extent does twin-informed AI-RAN orchestration improve continuity and worst-user QoE versus defined baselines?
\textbf{RQ3 (Multi-access resilience):} Under which outage, latency, and privacy conditions do terrestrial, offline, and NTN fallback policies help or fail?

\subsection{Contributions (methods only)}
\begin{itemize}
  \item Canonical cross-repository JSON Schema contracts and provenance for Edge-IO $\rightarrow$ twin $\rightarrow$ AI-RAN $\rightarrow$ NTN resilience paths.
  \item A preregistered statistical analysis plan with locked primary outcome \texttt{recovery\_time\_s} and baseline registries frozen before authentic pilot inspection.
  \item A 54-cell pilot design matrix with explicit consent, privacy scan, and session-mode rules separating calibration, rehearsal, and eligible pilot sessions.
  \item Honest claim boundaries: synthetic Gate~4 infrastructure results are labeled separately from pending Gate~3 field evidence.
\end{itemize}

Outcome tables and effect sizes are withheld until authentic Gate~3 pilot data exist (\textbf{RESULTS\_PENDING\_AUTHENTIC\_GATE3\_DATA}).

\section{Related Work}
\label{sec:related}

\subsection{IMT-2030 and 6G vision}
ITU-R Recommendation M.2160 defines the framework and overall objectives for IMT-2030 development~\cite{itu_m2160}.
Survey literature outlines AI-native air interfaces, extreme connectivity, and integrated sensing and communication as research directions rather than deployed capabilities~\cite{saad2020vision,strinati20216g}.
We use \emph{6G-aligned} and \emph{IMT-2030-aligned} wording only.

\subsection{Open RAN and AI-RAN experimentation}
The O-RAN Alliance architecture disaggregates the RAN into open interfaces suitable for policy and ML experimentation~\cite{oran_arch_v5}.
This artifact implements \emph{AI-assisted O-RAN-style} control experiments in simulation; it does not claim near-real-time RIC deployment on commercial networks.

\subsection{Edge computing and measurement}
ETSI MEC defines a framework for edge-hosted applications and orchestration~\cite{etsi_mec003}.
Companion work in \texttt{edge-io-measurement-node} specifies opt-in, schema-validated telemetry without payload capture; we reference that schema rather than redefining it here.

\subsection{Non-terrestrial networks}
3GPP TR~38.821 documents NR solutions for NTN scenarios including satellite and aerial platforms~\cite{3gpp_tr38821}.
Our resilience baselines (\texttt{terrestrial\_only}, \texttt{always\_ntn\_on\_terrestrial\_failure}, etc.) are evaluated in source-validated simulation paths; live NTN control is not claimed.

\subsection{Reproducibility and preregistration}
Preregistered outcomes, holdout registries, and negative-result preservation follow open-science reproducibility practices documented in \texttt{evaluation/EVALUATION\_PREREGISTRATION.md} (formal preregistration citation: CITATION\_NEEDED).
This repository locks \texttt{PRIMARY\_OUTCOME\_LOCK.json} and \texttt{EVALUATION\_PREREGISTRATION.md} before complete pilot result inspection.

\section{System Architecture and Contracts}
\label{sec:architecture}

\subsection{Umbrella orchestration role}
\texttt{gunnchos-7gc-ai-ran-field-kit} is the canonical contract authority linking:
\begin{itemize}
  \item \texttt{edge-io-measurement-node} --- privacy-preserving device telemetry (RQ1),
  \item \texttt{7gc-digital-twin} --- twin-state construction (RQ1--RQ2),
  \item \texttt{spectrumx-ai-ran-gary} --- AI-RAN policy baselines and ablations (RQ2--RQ3),
  \item \texttt{ntn-resilience-sim} --- multi-access resilience simulation (RQ3).
\end{itemize}
Optional PHY supporting studies (e.g., ReadyGary beam selection) may inform assumptions but are not required for thesis completion.

\subsection{Gate model}
\begin{table}[t]
\centering
\caption{Research gates and evidentiary meaning}
\label{tab:gates}
\begin{tabular}{ll}
\hline
\textbf{Gate} & \textbf{Meaning} \\
\hline
Gate 1 & Locked thesis, hypothesis, RQ1--RQ3 \\
Gate 2 & Executable schemas + integrated pipeline \\
Gate 3 & Authentic controlled-device pilot evidence \\
Gate 4 & Preregistered evaluation on eligible evidence \\
Gate 5/6 & Release archive, DOI, non-author reproduction \\
\hline
\end{tabular}
\end{table}

\subsection{Canonical schemas}
Contracts under \texttt{contracts/} include measurement session context, Edge-IO batches, twin bundles, AI-RAN decision bundles, and Gate status reports.
All integrated runs emit SHA-256 provenance via \texttt{scripts/verify\_provenance.py}.
Invalid fixtures under \texttt{fixtures/invalid/} document rejection reasons for contract tests.

\subsection{Integrated pipeline}
The default reproduction path executes:
\begin{verbatim}
make test && make integrated-pipeline \
  EDGE_INPUT=fixtures/valid/edge_measurement_batch.valid.json
\end{verbatim}
Outputs land under \texttt{results/integrated/} with evidence labels \texttt{synthetic} or \texttt{controlled\_device\_measurement}.
Synthetic outputs support Gate~4 infrastructure readiness only; they must not be presented as field pilot results.

\section{System Model and Notation}
\label{sec:system_model}

We formalize the artifact as a closed-loop measurement--twin--policy pipeline over discrete probe intervals.
Table~\ref{tab:notation} summarizes symbols used throughout the manuscript.
All empirical quantities derived from device probes are distinguished from model-estimated fields such as \texttt{expected\_recovery\_time\_s} in resilience decision bundles.

\begin{table}[t]
\centering
\caption{Core notation (methods manuscript)}
\label{tab:notation}
\small
\begin{tabular}{ll}
\hline
\textbf{Symbol} & \textbf{Meaning} \\
\hline
$\mathcal{U}$ & set of participant devices (clustered by day/zone) \\
$\mathcal{S}$ & service/workload profiles (\texttt{learn}, \texttt{create}, \texttt{sense}) \\
$\mathcal{M}$ & access modes: terrestrial, local edge, NTN, offline \\
$\mathbf{x}_t$ & Edge-IO probe at timestamp $t$ (latency, flags, intent) \\
$\mathcal{T}_t$ & digital-twin state bundle at $t$ (identifier-free) \\
$\pi$ & AI-RAN policy over RAN-style actions \\
$\rho$ & multi-access resilience policy over $\mathcal{M}$ \\
$\mathcal{C}$ & pilot matrix cell $(z,d,c,w)$: zone, day, condition, workload \\
\hline
\end{tabular}
\end{table}

\subsection{Session and matrix structure}
Each eligible pilot session $s$ is indexed by matrix cell $\mathcal{C}_s \in \mathcal{Z}\times\mathcal{D}\times\mathcal{N}\times\mathcal{S}$ with $|\mathcal{Z}|=3$ zones, $|\mathcal{D}|=3$ days, $|\mathcal{N}|=2$ network conditions, and $|\mathcal{S}|=3$ workload profiles (54 cells total).
Sessions with \texttt{session\_mode=PILOT} after consent, privacy scan, and assignment-hash verification contribute to Gate~3 evidence; calibration and rehearsal sessions are excluded from the matrix count.

\subsection{Probe timeline and availability}
For session $s$, ordered probes $\{\mathbf{x}_t\}_{t=1}^{T_s}$ are emitted by \texttt{edge-io-measurement-node} at workload-defined cadence (minimum 50 samples over 300~s per workload protocol).
Availability at probe $t$ is a Boolean function $a(\mathbf{x}_t)$:
\begin{equation}
a(\mathbf{x}_t)=
\begin{cases}
\texttt{service\_available}, & \text{if field present} \\
\neg\texttt{probe\_timeout}\land(\texttt{latency\_ms}\neq\emptyset), & \text{otherwise}.
\end{cases}
\end{equation}
Outage intervals begin at the first unavailable probe after availability and end at the first subsequent available probe, or continue until session end (right censoring).

\subsection{Primary and secondary outcomes}
The amended primary outcome is \textbf{total\_service\_outage\_time\_s}: the sum of outage-interval durations within session $s$ (lower is better).
Amendment~001 documents the change from the prior label \texttt{recovery\_time\_s}, which conflated total outage duration with a single recovery-time estimand; the amendment is internally valid and awaiting final author approval before physical collection.
The secondary outcome is \textbf{time\_to\_recovery\_s} at the event level: duration from outage onset to first available probe for \emph{completed} recoveries only.
Outages still active at the final probe are right-censored and excluded from completed-recovery summaries (no imputation as observed recoveries).
Physical computation follows \texttt{scripts/compute\_recovery\_time.py} on frozen \texttt{edge\_measurement\_batch} schema fields.

\subsection{Decision pipeline}
The integrated Gate~2 path maps $\mathbf{x}_t \rightarrow \mathcal{T}_t \rightarrow \pi(\mathcal{T}_t) \rightarrow \rho(\mathcal{T}_t,\pi)$, producing schema-validated \texttt{airan\_decision\_bundle} and \texttt{resilience\_decision\_bundle} artifacts with SHA-256 provenance.
AI-RAN policies $\pi$ select RAN-style actions (resource blocks, scheduling, compute placement); resilience policies $\rho$ select $\texttt{selected\_mode}\in\mathcal{M}$ subject to contract-recorded constraints and rejected-alternative audit trails.

\section{Orchestration Objectives, Constraints, and Multi-Access Logic}
\label{sec:orchestration}

\subsection{Service-continuity objective}
Twin-informed orchestration targets \emph{service continuity}: sustaining workload-appropriate connectivity and local-edge responsiveness under degraded conditions.
The AI-RAN layer optimizes predicted continuity-related metrics (\texttt{service\_continuity\_score}, scheduling and placement actions) subject to explicit constraint violations recorded in \texttt{constraint\_violations}.
The resilience layer selects among terrestrial, local-edge, NTN, and offline continuation modes to maximize \texttt{continuity\_score} while respecting energy, privacy, and feasibility gates encoded in \texttt{rejected\_alternatives}.
Model outputs such as \texttt{expected\_recovery\_time\_s} are \emph{not} substituted for physical \texttt{total\_service\_outage\_time\_s}; they support policy comparison and simulation only until Gate~3 evidence is frozen.

\subsection{Fairness, energy, reliability, and privacy constraints}
Constraints are enforced at decision time and audited in bundle schemas:
\begin{itemize}
  \item \textbf{Fairness:} AI-RAN bundles expose \texttt{fairness\_index} and ablation \texttt{fairness\_constraint} toggles worst-user protection versus uniform allocation; violations append to \texttt{constraint\_violations}.
  \item \textbf{Energy:} Resilience bundles record \texttt{estimated\_energy\_cost\_j}; ablation \texttt{energy\_constraint} removes the energy ceiling from candidate ranking.
  \item \textbf{Reliability:} Predicted reliability derives from loss models (\texttt{reliability\_from\_loss}: $1-\text{predicted\_packet\_loss\_pct}/100$) and informs mode rejection when below policy thresholds.
  \item \textbf{Privacy:} \texttt{estimated\_privacy\_cost} $\in[0,1]$ ranks alternatives; telemetry remains identifier-free per Edge-IO schema; ablation \texttt{remove\_privacy\_constraint\_simulated\_decision\_constraint\_only} is a simulated toggle---live privacy controls stay enabled during collection.
\end{itemize}
Ablation registry entries (\texttt{physical\_telemetry}, \texttt{digital\_twin\_context}, \texttt{continuity\_objective}, \texttt{local\_edge\_option}, \texttt{ntn\_option}, \texttt{workload\_intent}) isolate component contributions without inventing outcome numbers here.

\subsection{Terrestrial, edge, and NTN decision logic}
Resilience policy $\rho$ evaluates candidates in registry order for baseline \texttt{service\_aware\_multi\_access} (proposed path):
\begin{enumerate}
  \item \textbf{Terrestrial:} default when predicted latency and capacity satisfy workload minima and terrestrial link state is nominal.
  \item \textbf{Local edge (\texttt{local\_edge\_wifi}, \texttt{degraded\_local}):} considered when terrestrial probes indicate elevated latency or loss but a local-edge endpoint responds within \texttt{local\_edge\_response\_ms} budgets recorded in Edge-IO batches.
  \item \textbf{NTN fallback (\texttt{ntn\_fallback}):} eligible when terrestrial \emph{and} local-edge paths fail triggers (timeout, capacity floor, or policy-specific \texttt{fallback\_trigger}), and NTN assumptions in the bundle pass energy and privacy gates.
  \item \textbf{Offline continuation (\texttt{offline\_continuation}, \texttt{delayed\_synchronization}):} selected when no radio path meets continuity minima; retained versus suspended service features are enumerated for audit.
\end{enumerate}
Baselines \texttt{terrestrial\_only}, \texttt{terrestrial\_then\_offline}, \texttt{always\_ntn\_on\_terrestrial\_failure}, and \texttt{priority\_class\_fallback} instantiate strict subsets of this candidate set for controlled comparison.
\texttt{oracle\_hindsight\_analysis\_only} is analysis-only and never deployable.

\subsection{Policy pseudocode}
Algorithm~\ref{alg:resilience} summarizes the service-aware multi-access template implemented in \texttt{ntn-resilience-sim} (deterministic given inputs and policy name).

\begin{algorithm}[t]
\caption{Service-aware multi-access resilience selection (template)}
\label{alg:resilience}
\begin{algorithmic}[1]
\REQUIRE twin state $\mathcal{T}$, AI-RAN decision $d$, policy registry entry $P$
\ENSURE resilience bundle with \texttt{selected\_mode}, audit fields
\STATE $\mathcal{C} \leftarrow$ \textsc{BuildCandidates}($\mathcal{T}, d, P$) \COMMENT{terrestrial, edge, NTN, offline}
\FOR{each candidate $m \in \mathcal{C}$}
  \STATE $(\hat{L}_m, \hat{C}_m, E_m, R_m, \Pi_m) \leftarrow$ \textsc{EstimateMetrics}($m, \mathcal{T}$)
  \IF{\textsc{ViolatesFairness}($m$) \OR \textsc{ViolatesEnergy}($m, E_m$)}
    \STATE \textsc{Reject}($m$, reason); \textbf{continue}
  \ENDIF
  \IF{\textsc{ViolatesReliability}($m, R_m$) \OR \textsc{ViolatesPrivacy}($m, \Pi_m$)}
    \STATE \textsc{Reject}($m$, reason); \textbf{continue}
  \ENDIF
  \STATE score($m$) $\leftarrow$ \textsc{ContinuityScore}($m, \hat{L}_m, \hat{C}_m$, workload intent)
\ENDFOR
\STATE $m^\star \leftarrow \arg\max_{m \in \mathcal{C} \setminus \text{rejected}}$ score($m$)
\RETURN bundle(\texttt{selected\_mode}=$m^\star$, \texttt{rejected\_alternatives}, hashes)
\end{algorithmic}
\end{algorithm}

AI-RAN twin-informed policy $\pi_{\text{twin}}$ extends network-only scoring with twin-derived features (zone context, probe cadence, workload class) while sharing the same constraint audit structure; static and network-only baselines skip twin features; service-priority elevates continuity class weights; optimization-based runs constrained search.

\subsection{Computational complexity}
Let $|\mathcal{C}| \le 4$ access modes and $|\mathcal{A}|$ denote AI-RAN discrete action facets (resource blocks, scheduling tiers) bounded by fixture registry sizes.
Per integrated run, AI-RAN policy evaluation is $O(|\mathcal{A}|)$ under tabulated baseline implementations; resilience selection is $O(|\mathcal{C}|)$ with constant-factor metric estimation.
Twin construction is $O(T_s)$ in probe count for session $s$.
Gate~4 batch evaluation over $B$ baselines and $A$ ablations is $O(B \cdot A \cdot T_s)$ for metric recomputation-dominated workloads---tractable for the 54-cell pilot matrix on commodity hardware; decision \texttt{runtime\_s} and \texttt{peak\_memory\_mb} are logged per bundle for empirical complexity verification (no numeric results reported here).

\section{Measurement and Pilot Protocol}
\label{sec:measurement}

\subsection{Evidence tiers}
We distinguish smoke CI tests, synthetic reproducible benchmarks, controlled device measurements, and operational RAN claims (the last is \emph{not claimed}).

\subsection{Session modes}
Only sessions with \texttt{session\_mode=PILOT} count toward the 54-cell matrix after consent, privacy scan, and assignment-hash verification.
Calibration and rehearsal sessions support infrastructure validation and are explicitly excluded from pilot counts.

\subsection{Pilot matrix design}
The matrix spans:
\begin{itemize}
  \item 3 zones (\texttt{zone\_a}, \texttt{zone\_b}, \texttt{zone\_c}) with public labels pending human approval,
  \item 2 network conditions (\texttt{wifi\_normal}, \texttt{wifi\_degraded}) applied via scripted condition scripts,
  \item 3 collection days (\texttt{day\_01}--\texttt{day\_03}),
  \item 3 workload profiles (\texttt{learn}, \texttt{create}, \texttt{sense}), each 300~s.
\end{itemize}
Authoritative cell definitions live in \texttt{pilot/54\_CELL\_ASSIGNMENT\_MATRIX.csv}.
At manuscript freeze: \textbf{0/54} eligible cells collected; \textbf{HUMAN\_ACTION\_REQUIRED}.

\subsection{Workload protocols}
Each workload specifies periodic request cadence, timeout handling, and success criteria (duration $\geq$300~s, samples $\geq$50, explicit nulls for unavailable metrics).
Details are in \texttt{pilot/WORKLOAD\_PROTOCOLS.md}.

\subsection{Privacy and ethics}
Collection follows opt-in consent, no payload inspection, no direct identifiers, aggregated public reporting, and receive-only RF boundaries where SDR observation is discussed (see \texttt{docs/PRIVACY\_AND\_ETHICS.md}; formal IRB citation: CITATION\_NEEDED).
Raw device exports remain in gitignored \texttt{datasets/controlled/raw-private/}; sanitized derivatives require human approval before publication.

\subsection{Network condition protocol}
Condition application and restoration use \texttt{pilot/scripts/apply\_condition.sh}, \texttt{verify\_condition.sh}, and \texttt{restore\_condition.sh} with parameters in \texttt{pilot/configs/network\_conditions.yaml}.
Deviations are logged in \texttt{datasets/controlled/protocol-deviations/PROTOCOL\_DEVIATION\_LOG.md}.

\section{Evaluation Design and Statistical Plan}
\label{sec:evaluation}

\subsection{Primary claim and outcomes}
The preregistered primary claim states that twin-informed service-aware orchestration reduces service outage duration relative to static, network-only, and service-priority policies under defined degraded connectivity, subject to energy, fairness, privacy, and reliability constraints.
\textbf{Primary outcome:} \texttt{total\_service\_outage\_time\_s}---total seconds the service is classified unavailable within an eligible session (lower is better), computed from physical Edge-IO probes via \texttt{scripts/compute\_recovery\_time.py}.
\textbf{Secondary outcome:} \texttt{time\_to\_recovery\_s}---per-outage recovery duration for completed recoveries only.
Amendment~001 supersedes the prior primary label \texttt{recovery\_time\_s}, which summed unavailable intervals but was named like a single recovery-time estimand; the original lock remains byte-preserved under \texttt{evaluation/amendments/amendment\_001/} pending final author approval before collection.
Practical significance for the primary metric uses a 5~s minimal interesting difference (frozen with preregistration thresholds).

\subsection{Baseline rationale}
Baselines are registered before authentic result inspection to prevent post-hoc policy selection:
\begin{itemize}
  \item \texttt{static\_uniform} --- equal allocation ignoring twin and service intent (lower bound on adaptivity).
  \item \texttt{network\_only} --- reacts to network KPIs without twin context (isolates RQ1 contribution).
  \item \texttt{service\_priority} --- elevates continuity class without twin features (isolates intent-only path).
  \item \texttt{optimization\_based} --- constrained search baseline for RQ2 comparison.
  \item \texttt{twin\_informed} --- proposed AI-RAN policy using identifier-free twin state.
\end{itemize}
Resilience baselines span \texttt{terrestrial\_only} through \texttt{service\_aware\_multi\_access}, with \texttt{oracle\_hindsight\_analysis\_only} reserved for upper-bound analysis only.
Registries: \texttt{evaluation/BASELINE\_REGISTRY.yaml}, \texttt{evaluation/ABLATION\_REGISTRY.yaml}, \texttt{evaluation/HOLDOUT\_REGISTRY.yaml}.

\subsection{Holdout and ablation rationale}
Leave-one-out holdouts (\texttt{leave\_one\_day\_out}, \texttt{leave\_one\_zone\_out}, \texttt{leave\_one\_condition\_out}, \texttt{leave\_one\_workload\_out}) test generalization across matrix axes without treating 54 sessions as independent participants.
\texttt{stress\_scenario\_holdout} excludes configured stress scenarios from primary fitting (\texttt{evaluation/configs/held\_out\_scenarios.yaml}).
Ablations remove one component at a time (telemetry, twin, continuity objective, fairness, energy, local edge, NTN, workload intent) to attribute failures and neutral results.
Split leakage checks and duplicate detection are mandatory before any primary-outcome inference (\texttt{evaluation/LEAKAGE\_AND\_DUPLICATE\_REPORT.md}).

\subsection{Statistics, clustering, and censoring}
Analysis follows \texttt{evaluation/STATISTICAL\_ANALYSIS\_PLAN.md}:
\begin{itemize}
  \item Grouped summaries by day and zone with cluster-aware treatment of repeated measures.
  \item Bootstrap and repeated-seed uncertainty (95\% CI, 1000 resamples, seed~0 per \texttt{evaluation/configs/metrics.yaml}).
  \item Effect sizes and practical significance against \texttt{PRACTICAL\_SIGNIFICANCE\_THRESHOLDS.yaml}.
  \item Missing-data and protocol-deviation audits (\texttt{evaluation/MISSING\_DATA\_ANALYSIS.md}).
\end{itemize}
\textbf{Censoring treatment:} Outages ongoing at session end are right-censored.
Censored events contribute to \texttt{total\_service\_outage\_time\_s} (duration to last probe) but \emph{do not} supply observed \texttt{time\_to\_recovery\_s} values.
Summary statistics for the secondary metric exclude censored events; sensitivity analyses may report censoring rate and interval-censoring bounds ($\pm$ half median probe interval) without imputing completed recoveries.
Negative and neutral results are preserved per \texttt{evaluation/NEGATIVE\_AND\_NEUTRAL\_RESULTS.md}.

\subsection{Synthetic Gate~4 infrastructure}
Gate~4 automation demonstrates evaluation plumbing on synthetic or fixture-backed inputs.
Those runs are labeled \texttt{GATE4\_EVALUATION\_READY} infrastructure, not authentic pilot completion.
No numeric results from Gate~4 are imported into this methods manuscript.

\section{Results}
\label{sec:results}

\textbf{RESULTS\_PENDING\_AUTHENTIC\_GATE3\_DATA}

Authentic controlled-device pilot sessions required for primary-outcome analysis have not been collected at manuscript freeze.
The eligible pilot matrix status is \textbf{0/54} cells.
Calibration evidence (e.g., one valid Pixel~6a calibration session) supports infrastructure validation only and must not be used to claim pilot completion or causal superiority.

When authentic Gate~3 data exist, this section will report:
\begin{itemize}
  \item pilot coverage against \texttt{pilot/54\_CELL\_ASSIGNMENT\_MATRIX.csv},
  \item primary outcome \texttt{recovery\_time\_s} by baseline and holdout split,
  \item uncertainty intervals, effect sizes, and practical significance against locked thresholds,
  \item negative and neutral results per \texttt{evaluation/NEGATIVE\_AND\_NEUTRAL\_RESULTS.md},
  \item failure boundaries per \texttt{evaluation/FAILURE\_BOUNDARIES.md}.
\end{itemize}

No tables, figures, p-values, or effect sizes are reported here to avoid inventing empirical numbers.

\section{Limitations, Ethics, and Threats to Validity}
\label{sec:limitations}

\subsection{Scope limits}
This artifact is a research prototype, not commercial 6G, carrier-grade AI-RAN, or citywide deployment infrastructure.
Gary/NWI anchor context supports digital-equity motivation but does not justify global generalization.

\subsection{Evidence limits}
Synthetic benchmarks and Gate~4 infrastructure runs demonstrate reproducible engineering and evaluation plumbing only.
Until Gate~3 authentic pilot data exist, all outcome claims remain \textbf{RESULTS\_PENDING\_AUTHENTIC\_GATE3\_DATA}.

\subsection{Design threats}
Repeated measures within days and zones induce clustering; the statistical plan addresses this but cannot remove locale-specific confounds.
Network condition scripts approximate degradation; they do not emulate full RF propagation diversity.
NTN resilience paths are simulation-backed unless and until validated against controlled NTN-capable environments.
Interval censoring from finite probe cadence bounds temporal precision of outage onsets and recoveries.

\subsection{Privacy, consent, and human subjects}
Structured field campaigns require IRB or equivalent approval before identifiable participant collection (CITATION\_NEEDED).
Consent versions, opt-in boundaries, and deletion/retention policies are documented in \texttt{docs/PRIVACY\_AND\_ETHICS.md}.
Public release depends on sanitized derivatives and consent-summary approval.
Missing cells, protocol deviations, and consent withdrawals must be reported rather than imputed without documented sensitivity analysis.
Simulated privacy-constraint ablations never disable live telemetry protections.

\subsection{Reproducibility limits}
Clean-checkout author reproduction, non-author reproduction, immutable release tarball, and DOI deposit remain Gate~5/6 requirements and are pending at manuscript freeze.
PDF builds use pinned tooling (\texttt{scripts/build\_paper.sh}) with \texttt{SOURCE\_DATE\_EPOCH=0} for reproducibility where supported; residual PDF byte differences may arise from hyperref object IDs or tool-specific metadata (documented in \texttt{PAPER\_BUILD\_REPORT.md}).

\section{Conclusion}
\label{sec:conclusion}

We presented a methods-ready, honesty-bounded artifact for resilience-aware, human-centric AI-RAN orchestration research spanning privacy-preserving measurement, twin-informed policy evaluation, and NTN fallback simulation.
Gate~1 thesis elements, Gate~2 contracts, Gate~3 pilot protocol, and Gate~4 preregistration are in place; empirical pilot outcomes remain \textbf{RESULTS\_PENDING\_AUTHENTIC\_GATE3\_DATA}.
The amended primary outcome \texttt{total\_service\_outage\_time\_s} with secondary \texttt{time\_to\_recovery\_s} aligns physical probe semantics with preregistered claims (Amendment~001, approval pending).

Future work completes the 54-cell matrix under human-approved zones and dates, executes the locked statistical analysis plan on eligible evidence only, and pursues release/archive steps with independent reproduction before outcome claims appear in a results manuscript.


\bibliographystyle{IEEEtran}
\bibliography{references}

\appendices
\section{Reproducibility Appendix}
\label{app:repro}

\subsection{Environment}
See repository \texttt{ENVIRONMENT.md} for Python version, dependency pins, and sibling repository layout under \texttt{REPOS\_ROOT}.
Integrated runs require sibling repos \texttt{edge-io-measurement-node}, \texttt{7gc-digital-twin}, \texttt{spectrumx-ai-ran-gary}, and \texttt{ntn-resilience-sim} at commits recorded in \texttt{PROVENANCE\_LOCK.json}.

\subsection{Author reproduction (pending)}
\texttt{AUTHOR\_REPRODUCTION\_REPORT.md} records a clean-checkout run outcome; status at freeze: \textbf{PENDING}.

\subsection{Non-author reproduction (pending)}
\texttt{NON\_AUTHOR\_REPRODUCTION\_REPORT.md} is an authentic empty pending template until an independent replicator files results.

\subsection{Core commands}
\begin{verbatim}
pip install -r requirements.txt
make test
make contract-test
make integrated-pipeline \
  EDGE_INPUT=fixtures/valid/edge_measurement_batch.valid.json
make gate1-validate
python3 scripts/compute_recovery_time.py \
  --input fixtures/valid/edge_measurement_batch.valid.json
\end{verbatim}

\subsection{Outcome metric reproduction}
Physical outcomes are recomputed from Edge-IO batches only.
The CLI above emits \texttt{total\_service\_outage\_time\_s}, event-level \texttt{time\_to\_recovery\_s}, censoring flags, and probe cadence diagnostics without referencing model \texttt{expected\_recovery\_time\_s} fields.

\subsection{Pilot collection (human-gated)}
Pilot counting requires \texttt{pilotctl.py} workflows, approved zone/date tokens, and storage under \texttt{datasets/controlled/raw-private/} (gitignored).
Public artifacts document protocol only until sanitized publication approval.

\subsection{Paper build reproducibility}
PDF build uses pinned Tectonic (\texttt{scripts/build\_paper.sh}, version pinned in-script) or a digest-pinned TeX Live Docker image; local \texttt{pdflatex}/\texttt{bibtex} remains an optional fallback.
Builds export \texttt{SOURCE\_DATE\_EPOCH=0} (and \texttt{TZ=UTC}) before compilation to stabilize timestamps where tools honor them.
Unavoidable byte-level PDF differences may still occur from hyperref anchor generation, filesystem XMP metadata, or Tectonic cache state; \texttt{PAPER\_BUILD\_REPORT.md} records page count (required 6--8 inclusive), SHA-256, and build method each run.
See \texttt{paper/Makefile} and \texttt{paper/README.md}.


\end{document}
